\subsubsection{比較経路としての圧縮--拡散型再初期化と厚み補正}
\label{sec:dgr_thickness}
% ─────────────────────────────────────────────────────────────────────────────

圧縮--拡散型再初期化と Direct Geometric Reinitialization
（DGR）は，本稿の標準段階 D/F 経路ではない．
これらは，圧縮項と拡散項を逐次適用する比較経路が
界面幅をどのように変えるかを調べるための参照であり，
標準計算へ暗黙の後処理として加えない．

DGR を比較経路として明示する場合だけ，界面帯の有効幅を
\begin{equation}
  \varepsilon_{\mathrm{eff}}
  =
  \operatorname{median}_{0.05<\psi<0.95}
  \frac{\psi(1-\psi)}{|\bnabla\psi|}
  \label{eq:dgr_effective_width}
\end{equation}
で評価し，
\begin{equation}
  \phi_{\mathrm{sdf}}
  =
  \varepsilon\ln\frac{1-\psi}{\psi}
  \frac{\varepsilon_{\mathrm{eff}}}{\varepsilon},
  \qquad
  \psi_{\mathrm{new}}=H_\varepsilon(-\phi_{\mathrm{sdf}})
  \label{eq:dgr_width_rescale}
\end{equation}
という幅再調整として定義する．
この定義は，入力が
$\psi=H_{\varepsilon_{\mathrm{eff}}}(-\phi_{\mathrm{true}})$
という一様幅のプロファイルである場合に限り，
$H_\varepsilon(-\phi_{\mathrm{true}})$ を復元する．
狭い形状特徴，局所 fold，表面張力で振動する界面では
単一の $\varepsilon_{\mathrm{eff}}$ が幾何を代表しないため，
DGR を標準閉包として採用しない．

標準アルゴリズムでは，段階 D が零点集合から符号距離場を再構成し，
段階 F が式~\eqref{eq:cls_phi_mass_correction} の
$\phi$ 空間一様シフトで液相体積を閉じる．
この節は，圧縮--拡散型比較経路と標準 Ridge--Eikonal 経路を
混同しないための境界条件を与えるために残す．


% ─────────────────────────────────────────────────────────────────────────────
\subsubsection{Ridge--Eikonal 再初期化：Eikonal 基盤と距離再構成の統合定式化}
\label{sec:xi_sdf_reinit}
\phantomsection\label{sec:eikonal_reinit}% 外部参照保護
% ─────────────────────────────────────────────────────────────────────────────

\noindent\textbf{前章までの連続定式化との関係：}
本節は §~\ref{sec:reinit}（§3.3 再初期化方程式）と
§~\ref{sec:ridge_eikonal}（§3.4 Ridge--Eikonal 連続定式化）の離散定式化である．
本稿の補助距離再構成はここで定式化する Ridge--Eikonal 系
（前段の Eikonal 距離構築 + 零点集合からの $\psi$ 再構成）であり，
DCCD/圧縮--拡散型の比較経路は §~\ref{sec:dissipative_ccd}（§4.4）に集約する．

\paragraph*{Eikonal 基盤と継承される理論保証}
Ridge--Eikonal は Eikonal PDE による符号距離関数構築を前段に持つ．
本項では Ridge--Eikonal が継承する Eikonal 法の基盤と理論保証を整理する．
Eikonal 方程式 \cite{Sussman1994}
\begin{equation}
  \frac{\partial\phi}{\partial\tau} + \operatorname{sgn}(\phi_0)\bigl(|\bnabla\phi| - 1\bigr) = 0
  \label{eq:eikonal_pde}
\end{equation}
は零点集合（界面位置）を不変に保ちつつ $|\bnabla\phi| \to 1$ を回復する（符号距離関数化）．
これによりフォールド（$|\bnabla\phi|\to 0$ 領域）も自然に修正される．

\begin{tcolorbox}[algbox, title={Ridge--Eikonal 距離再構成の基本手順}]
\begin{enumerate}
  \item \textbf{ロジット反転：}$\phi = \varepsilon\ln((1-\psi)/\psi)$
        （値域制限は数値特異性を避けるための入力整形であり，
        零点集合は後続の交差検出から決める）
  \item \textbf{周期商空間への射影：}周期軸集合を $P$ とすると，
        離散スカラーは
        \begin{equation}
          \mathcal{Q}_P
          =
          \{q_h\mid
          q_{i_1,\ldots,N_a,\ldots,i_d}
          =
          q_{i_1,\ldots,0,\ldots,i_d}\quad (a\in P)\}
          \label{eq:periodic_quotient_space}
        \end{equation}
        に属する．Eikonal/FMM 入力の $\phi$ と再構成後の $\psi$ は
        各呼び出しで $\Pi_P q_h\in\mathcal{Q}_P$ へ射影し，
        周期像ノードを独立自由度として扱わない．
  \item \textbf{距離場構築：}
        一様格子では §~\ref{sec:xi_sdf_reinit} の ξ-SDF，
        非一様格子や壁面接触を含む場合は
        §~\ref{sec:ridge_eikonal_nu_grid} の FMM/FSM 型距離再構成で
        零点集合から $\phi$ を構成する．
        Godunov/WENO--HJ 型の仮想時間反復は比較経路であり，
        標準段階 D の前提にはしない．
  \item \textbf{セル局所 $\varepsilon$ による再構成：}
        $\psi(i,j) = H_{\varepsilon(i,j)}(-\phi)$，
        $\varepsilon(i,j) = \varepsilon_\xi\max(h_x(i),\,h_y(j))$

        ここで $\varepsilon_\xi$ は §\ref{sec:cls_advection}
        の界面幅係数 $C_\varepsilon$ と同一のスカラー定数
        （$\varepsilon = C_\varepsilon\Delta x_{\min}$ より $\varepsilon_\xi = C_\varepsilon \approx 1.0\sim 2.0$）；
        本節では ξ 計算座標で $\varepsilon_\xi$ 表記を採用し，
        §\ref{sec:cls_advection}/§\ref{sec:intro} と表記体系のみ異なる
        同一物理量を表す．
        一様格子では $\varepsilon(i,j) = \varepsilon$（スカラーに帰着）；
        非一様格子では各セルの物理格子幅に比例した界面幅を実現する．
  \item \textbf{$\phi$ 空間質量補正：}
        $\delta\phi = -\Delta M / \int H'_{\varepsilon(\bx)}\,dV$，
        $\phi \leftarrow \phi + \delta\phi$，
        $\psi \leftarrow H_{\varepsilon(\bx)}(-\phi)$．
        ただし壁面接触点が存在する場合は
        $\phi \leftarrow \phi+\delta\phi\,\chi_{\mathrm{free}}$ とし，
        接触線帯では $\chi_{\mathrm{free}}=0$ として零集合を固定する
        （§~\ref{sec:ridge_eikonal_nu_grid} の壁閉包補正 W）．
\end{enumerate}
\end{tcolorbox}

\noindent\textbf{理論的保証：}
\begin{enumerate}
  \item \textbf{零点集合保存：}
        標準経路では，零点集合を入力の交差集合または物理壁接触点から定義し，
        距離構築はその集合を移動させない写像として扱う．
  \item \textbf{距離性：}
        受理された $\phi$ は $|\bnabla\phi|\approx1$ を満たし，
        再構成後の界面幅は局所幅 $\varepsilon(\bx)$ で決まる．
  \item \textbf{質量閉鎖（線形化補正）：}手順 4 の $\phi$ 空間補正
        $\delta\phi=-\Delta M/\!\int H'_{\varepsilon(\bx)}\,dV$ は
        $M(\phi)=\int H_{\varepsilon(\bx)}(-\phi)\,dV$ に対する 1 回の
        Newton 型線形化補正である．したがって小さい残差では
        $\sum\psi_{\mathrm{corr}}\approx\sum\psi_{\mathrm{old}}$ となるが，
        厳密閉鎖を要求する場合は残差を監視し，必要に応じて同じ
        $\phi$ 空間補正を反復する．壁面接触界面では補正重みを
        接触線帯で零にし，質量補正が接触点を移動しない
        制約付き線形化として用いる．
  \item \textbf{周期商空間閉性と対称性：}
        $\Omega=S^1\times[0,L_y]$ のような混合境界では，
        周期軸上の端点 $i=N_x$ は $i=0$ の像であり，
        離散再初期化作用素 $\mathcal{R}_h$ は
        \begin{equation}
          \mathcal{R}_h(\mathcal{Q}_P)\subset\mathcal{Q}_P,
          \qquad
          \mathcal{R}_h\Pi_P=\Pi_P\mathcal{R}_h
          \label{eq:reinit_periodic_quotient_closure}
        \end{equation}
        を満たさなければならない．さらに初期条件・格子・壁条件が
        半周期並進 $T_{N_x/2}$ または反射 $R_x$ に不変なら，
        $\mathcal{R}_h T_{N_x/2}=T_{N_x/2}\mathcal{R}_h$，
        $\mathcal{R}_h R_x=R_x\mathcal{R}_h$ が必要である．
        したがって FMM の Dirichlet 種点は実際の零交差集合
        $S_h(\phi_0)$ と物理壁接触拘束に限られ，補助リッジの局所極大を
        ``近いから'' という理由で $\phi=0$ 種点に昇格してはならない．
        これは Eikonal の零点集合保存性と周期商空間の同値関係を同時に保つための
        離散幾何条件である．
\end{enumerate}

\paragraph*{Eikonal 単独運用の限界：Godunov 由来の零セットドリフト}
Eikonal 法は $\sigma{=}0$ の受動移流・静的形状問題（Zalesak スロット等）では
体積保存・界面幅・形状復元のすべての条件を満たす有力手法であるが，
$\sigma>0$ で界面振動を伴う系では Godunov 差分の離散誤差により
零点集合が各呼び出しでわずかに移動し，反復累積で形状誤差を生む．
これは圧縮--拡散型比較経路の幅再調整誤差とは別の，
反復型距離化に固有の系統誤差である．
本制約は $\phi$ の局所操作を反復する全手法に共通であり，
\textbf{反復を一切伴わない解析的距離構築}を要請する——これが
Ridge--Eikonal（次項 ξ 空間符号距離関数法）の動機である．


\paragraph*{反復なし基底解：ξ 空間符号距離関数法}
反復累積誤差を原理的に排除するため，
本経路を選ぶ場合は，再初期化呼び出しごとに
現在の $\psi$ 場から零点交差を直接検出し，
格子インデックス空間（$\xi$ 空間）における符号距離を
解析的に構築する \textbf{ξ 空間符号距離関数（ξ-SDF）法}を採用する．
本手法は Ridge--Eikonal 系の一様格子基底形であり，反復を行わないため
累積ドリフトが構造的に存在しない．ガウス補助場 $\xiridge$ と非一様 FMM を用いる
完全形は §~\ref{sec:ridge_eikonal_nu_grid} で扱う．

\begin{tcolorbox}[algbox, title={ξ-SDF 再初期化アルゴリズム}]
\begin{enumerate}
  \item \textbf{ロジット反転：}
        $\phi_{\mathrm{raw}}(i,j) = \varepsilon\ln\!\bigl((1{-}\psi(i,j))/\psi(i,j)\bigr)$
        （飽和制限付き）
  \item \textbf{零点交差の検出（線形補間）：}$\phi_{\mathrm{raw}}$ の符号が変化する隣接セル対を探索し，
        格子座標上の交差位置を線形補間で決定する．
        $x$-方向（$(i,j)\to(i{+}1,j)$）の交差 $\xi$-座標：
        \begin{equation}
          \xi^* = i + \frac{|\phi_{\mathrm{raw}}(i,j)|}{|\phi_{\mathrm{raw}}(i,j)| + |\phi_{\mathrm{raw}}(i{+}1,j)|}
          \label{eq:xi_crossing_x}
        \end{equation}
        $y$-方向（$(i,j)\to(i,j{+}1)$）の交差 $\eta$-座標：
        \begin{equation}
          \eta^* = j + \frac{|\phi_{\mathrm{raw}}(i,j)|}{|\phi_{\mathrm{raw}}(i,j)| + |\phi_{\mathrm{raw}}(i,j{+}1)|}
          \label{eq:xi_crossing_y}
        \end{equation}
        全交差位置の集合を $S_\xi = \{(\xi^*_k,\,\eta^*_k)\}_{k=1}^{N_c}$ とする．
  \item \textbf{ξ距離場の構築：}各セル $(i,j)$ から最近傍零点交差までの
        ξ空間ユークリッド距離を計算し，$\phi_{\mathrm{raw}}$ の符号を付与する：
        \begin{equation}
          \phi_\xi(i,j) = \operatorname{sgn}\!\bigl(\phi_{\mathrm{raw}}(i,j)\bigr)
          \cdot \min_{(\xi^*,\eta^*)\in S_\xi}
          \sqrt{(i-\xi^*)^2 + (j-\eta^*)^2}
          \label{eq:xi_sdf}
        \end{equation}
  \item \textbf{$\psi$ 再構成：}定数 $\varepsilon_\xi = \varepsilon/h_{\min}$
        （= $C_\varepsilon$，§\ref{sec:cls_advection} 参照）を用いて
        \begin{equation}
          \psi_{\mathrm{new}}(i,j) = H_{\varepsilon_\xi}\!\bigl(-\phi_\xi(i,j)\bigr)
          = \frac{1}{1+\exp\!\bigl(\phi_\xi(i,j)/\varepsilon_\xi\bigr)}
          \label{eq:xi_sdf_psi}
        \end{equation}
  \item \textbf{$\phi_\xi$-空間質量補正（必要時）：}
        $\delta\phi = -\Delta M / \int H'_{\varepsilon_\xi}\,dV$，
        $\phi_\xi \leftarrow \phi_\xi + \delta\phi$，
        $\psi \leftarrow H_{\varepsilon_\xi}(-\phi_\xi)$．
        この最終補正は界面を一様に $\delta\phi$ だけ移動するため，
        後述の零交差保存命題は補正前の距離場に対する性質である．
\end{enumerate}
\end{tcolorbox}

\paragraph*{理論的性質}

\begin{description}
\item[命題~1（補正前距離場における線形補間零交差の保存）]
式~\eqref{eq:xi_crossing_x}--\eqref{eq:xi_crossing_y} で得た
零交差集合 $S_\xi$ 上では，質量補正前の距離場について
$\phi_\xi=0$ であり，
式~\eqref{eq:xi_sdf_psi} により $\psi_{\mathrm{new}}=0.5$ となる．

\textbf{前提}：$S_\xi$ は隣接セル対の符号反転から構成する．格子点で偶然
$\phi_{\mathrm{raw}}=0$ となる退化点は離散上の符号制限で片側セルへ帰属させ，
界面位置は隣接符号反転の線形補間として扱う．したがって保存対象は
連続界面そのものではなく，離散入力から構成した線形補間零交差集合である．

\textbf{証明：}任意の $(\xi^*,\eta^*)\in S_\xi$ に対して，
式~\eqref{eq:xi_sdf} の距離最小値は 0 である．よって
$\phi_\xi(\xi^*,\eta^*)=0$，式~\eqref{eq:xi_sdf_psi} より
$\psi_{\mathrm{new}}=1/(1+e^0)=0.5$．格子点上での界面位置は
この零交差集合を線形補間で再標本化したものとなる．
最終の $\phi$ 空間質量補正を適用する場合，零点集合は
$\phi_\xi=-\delta\phi$ へ一様に移動するため，保存対象は
補正前の $S_\xi$ と補正後の一様シフト量 $\delta\phi$ の組である．\qed

\item[命題~2（ξ-SDF 性質：$|\bnabla_\xi\phi_\xi|=1$）]
式~\eqref{eq:xi_sdf} で定義される $\phi_\xi$ は，格子インデックス空間において
$|\bnabla_\xi\phi_\xi|=1$ を満たす真の符号距離関数である（Voronoi 境界を除く）．

\textbf{証明：}$d(i,j;S_\xi) = \min_k\sqrt{(i-\xi^*_k)^2+(j-\eta^*_k)^2}$ は
$\xi$ 空間の点集合 $S_\xi$ への Euclid 距離場であり，
$|\bnabla d|=1$ はほぼすべての点で成立する（距離場の基本性質）．
$\phi_\xi = \operatorname{sgn}(\phi_{\mathrm{raw}}) \cdot d$ であり，$\operatorname{sgn}$ の符号変化は
$S_\xi$ 上（$d=0$）でのみ起きるため，$|\bnabla_\xi\phi_\xi|=|\bnabla d|=1$．\qed

\item[命題~3（無反復 $\Rightarrow$ ドリフト累積なし）]
ξ-SDF は反復仮想時間積分を行わない．
補正前の零交差集合は線形補間誤差 $\Ord{h^2}$ の範囲で
入力の符号反転集合から構成され，Godunov 差分の
$O(\Delta\tau\cdot n_{\mathrm{iter}})$ 累積ドリフトを生じない．
補正後の界面移動は段階 F の意図的な一様シフト
$\delta\phi$ として別に監視する．
\end{description}

\paragraph*{設計上の運用方針と $\varepsilon$ 幅効果}
ξ-SDF は反復累積誤差を持たない一方，正確な幅 $\varepsilon$ の界面を生成するため，
$\sigma>0$ で界面振動を伴う系では表面張力の空間集中
（$\sim\varepsilon^{-1}$ スケール）が PPE 投影残差を介して
体積誤差に伝播し得る．
この効果を標準経路の暗黙の正則化として処理してはならない．
幅拡大係数 $f$ は，Ridge--Eikonal 距離再構成の感度を測る
明示的な比較パラメタとしてのみ導入し，
\begin{equation}
  \psi_{\mathrm{new}}(i,j) = H_{f\varepsilon_\xi}\!\bigl(-\phi_\xi(i,j)\bigr)
  \label{eq:xi_sdf_wide}
\end{equation}
と書く．標準の幅保存経路は $f{=}1$ であり，$f>1$ を使う場合は
§~\ref{sec:val_capillary} の感度検証条件として明示し，標準計算と混同しない．
非一様格子では式~\eqref{eq:xi_sdf_wide} の定数幅をそのまま用いず，
§~\ref{sec:ridge_eikonal_nu_grid} の D3 非一様 FMM と
D4 空間依存 $\varepsilon_\text{local}(\bx)$ へ置き換える．

\paragraph*{適用範囲（設計指針）}
\begin{itemize}
  \item $\sigma=0$ 問題（受動移流・静的形状）：ξ-SDF（$f{=}1$）または FMM~\cite{Sethian1996}．
        反復なしで $C^0$ 連続な符号距離場を解析的に構築し，
        Godunov 由来のドリフト累積を排除する．
  \item $\sigma>0$ 問題（界面振動を伴う系）：Ridge--Eikonal/FMM/ξ-SDF の
        零点集合保存と段階 F の体積閉鎖を標準条件とする．
        $f>1$ の幅拡大や分裂法単独は比較・感度検証として明示し，
        標準経路の安定化として黙示的に採用しない．
\end{itemize}

% ─────────────────────────────────────────────────────────────────────────────
\subsubsection{再初期化の受理条件}
% ─────────────────────────────────────────────────────────────────────────────

\begin{tcolorbox}[mybox, title={受理条件：段階 D/F の受理前提}]
Ridge--Eikonal 再構成を受理する条件は，
\begin{enumerate}[noitemsep]
  \item 入力の零点集合または壁面接触点と整合すること，
  \item 再構成後の $\phi$ が距離性 $|\bnabla\phi|\approx1$ を満たすこと，
  \item 再構成後の閾値化相体積 $V_{\Gamma,h}$ が
        式~\eqref{eq:ridge_eikonal_sharp_volume_shift} の
        一様 SDF シフトで許容残差内へ閉じること，
  \item 再構成後の拡散 CLS 質量 $M_{\psi,h}$ が
        式~\eqref{eq:ridge_eikonal_diffuse_mass_width} の
        プロファイル幅補正で許容残差内へ閉じること，
  \item $Q_\Gamma/Q_{\Gamma,\mathrm{ref}}$ が採用した実行条件に対して
        許容範囲にあること，
  \item 曲率評価に使う界面帯が単調な $H_\varepsilon(-\phi)$ 型プロファイルとして
        保たれていること，
\end{enumerate}
である．反復型 Eikonal や圧縮--拡散型再初期化を比較経路として選ぶ場合も，
同じ受理量で標準 Ridge--Eikonal 経路と比較する．
\end{tcolorbox}

% ─────────────────────────────────────────────────────────────────────────────
% 時間積分関連 2 節は §7 に集約．本ファイルでは §7 への redirect 1 行のみ残す．
% ─────────────────────────────────────────────────────────────────────────────
\medskip
\noindent\textbf{再初期化の時間刻み制約：}
完全陽解法として再初期化方程式を仮想時間積分する場合の $\Delta\tau$ 安定性
（双曲 CFL + 拡散 CFL）と，物理 $\Delta t$ の対流 CFL・界面張力 CFL・粘性 CFL
（律速合成式）は時間積分章にまとめている．
仮想時間 $\Delta\tau$ の参考安定条件は \S~\ref{warn:cls_dtau_stability_v7}（§7）に，
物理 $\Delta t$ の律速合成式は \S~\ref{sec:dt_synthesis}（§7.9）に集約．
本稿の Ridge--Eikonal 補助経路（\S~\ref{sec:xi_sdf_reinit}）は
反復不要であり，仮想時間 CFL は物理 $\Delta t$ には乗らない．
